%=============================================================================
% THIS IS chapter{Gain amplification verification tool}
%
% May 2026 - created for STL integration
%=============================================================================
\chapter{Gain amplification verification tool}
%=============================================================================

\section{Introduction}

The \texttt{gain\_chk} tool compares a reference signal and a processed
signal and reports short-time gain differences per frequency band. The
tool is intended for objective verification of spectral gain behavior in
speech and audio processing chains.

Both inputs are raw 16-bit PCM files. The processing is frame-based with
20 ms frames and supports sampling frequencies 8, 16, 32 and 48 kHz.
Optionally, a VAD file can be provided to separate statistics for active
and inactive frames.

\section{Description of the algorithm}

\subsection{Frame and spectrum analysis}

Let $x_1[n]$ be the reference frame and $x_2[n]$ be the processed frame,
with frame length $L=\frac{F_s}{50}$ samples (20 ms). Each frame is
windowed with a Hamming window
\begin{equation}
w[n] = 0.54 - 0.46\cos\left(\frac{2\pi n}{L-1}\right), \quad n=0,\ldots,L-1.
\end{equation}

The tool computes the real FFT of each windowed frame after zero-padding
to the next power-of-two length. For each frequency bin $k$, the spectral
energy is
\begin{equation}
P_i[k] = \Re\{X_i[k]\}^2 + \Im\{X_i[k]\}^2, \quad i \in \{1,2\}.
\end{equation}

\subsection{Band energy accumulation}

The analysis bands are defined by upper frequency limits in the
\texttt{bands[]} table. Values in $(0,1]$ are interpreted as normalized
fractions of Nyquist ($F_s/2$). For each band $b$, the tool accumulates
average bin energy:
\begin{equation}
E_i^{(b)} = \frac{1}{N_b}\sum_{k \in \mathcal{B}_b} P_i[k],
\end{equation}
where $\mathcal{B}_b$ is the set of bins inside the band and $N_b$ its
number of bins.

For each frame class (all/active/inactive), band energies are accumulated
across frames:
\begin{equation}
\bar{E}_i^{(b)} = \sum_{\mbox{\scriptsize frames}} E_i^{(b)}.
\end{equation}

\subsection{Reported gain measures}

For each band, the gain ratio is
\begin{equation}
R^{(b)} = \frac{\bar{E}_2^{(b)} + E_{\min}}{\bar{E}_1^{(b)} + E_{\min}},
\end{equation}
where $E_{\min}$ is a small offset used for numerical robustness.

The reported outputs are:
\begin{equation}
G_{\mathrm{dB}}^{(b)} = 10\log_{10}\left(R^{(b)}\right),
\end{equation}
\begin{equation}
G_{\%}^{(b)} = \left(R^{(b)} - 1\right)\cdot 100.
\end{equation}

The same ratio is also computed on the total energy over all bands.

\subsection{VAD-conditioned processing}

When a VAD file is present (\texttt{-v}), one 16-bit flag is read per
frame. Frames with non-zero VAD are treated as active. Frames with zero
VAD are treated as inactive.

For inactive frames, a low-energy gate is applied on the reference frame:
if mean time-domain energy is below a fixed threshold, that frame is not
included in spectral statistics. This avoids unstable ratios in near-silent
segments.

\subsection{Threshold check mode}

When threshold options are enabled, two checks are available:
\begin{itemize}
\item \texttt{-a}: maximum allowed positive gain (amplification) in dB for active frames;
\item \texttt{-s}: maximum allowed attenuation in dB for inactive frames (requires VAD).
\end{itemize}

The tool returns a numeric code intended for automation scripts:
\begin{itemize}
\item 1: active fail and inactive fail;
\item 2: active pass and inactive fail;
\item 3: active fail and inactive pass;
\item 4: active pass and inactive pass.
\end{itemize}

\section{Implementation}

\subsection{Data and file format}

The tool expects raw header-less little or big endian 16-bit PCM files in
native host byte order, consistent with other STL command-line tools.

Mandatory input/output arguments:
\begin{itemize}
\item \texttt{-i} reference PCM file;
\item \texttt{-o} processed PCM file;
\item \texttt{-t} output text file;
\item \texttt{-r} sampling frequency in Hz.
\end{itemize}

Optional arguments:
\begin{itemize}
\item \texttt{-v} VAD file (one 16-bit flag per frame);
\item \texttt{-a} active-frame amplification threshold (dB);
\item \texttt{-s} inactive-frame attenuation threshold (dB).
\end{itemize}

\subsection{Command line syntax}

{\tt {\small
\begin{tabular}{llll}
gain\_chk & -i ref.pcm & -o proc.pcm & -t results.txt \\
          & -r Fs      & [-v vad.bin] & [-a A\_dB] [-s S\_dB] \\
\end{tabular}
}}

\subsection{Processing flow}

Pseudo code:
{\tt\small
\begin{verbatim}
Read and validate command line
Initialize frame length, bands, and Hamming window
for each complete 20 ms frame in both input files:
    read reference and processed PCM frame
    read VAD flag if provided (else active by default)
    optionally skip low-energy inactive frame
    apply Hamming window and zero padding
    compute FFT of both frames
    for each frequency band:
        accumulate average spectral energy in reference/processed
After loop:
    print per-band and total gain in dB and percent
    optionally evaluate thresholds and return status code
\end{verbatim}
}

\section{Tests and portability}

The STL CMake integration defines regression tests for \texttt{gain\_chk}
using bundled PCM material. Tests run at 48 kHz and compare expected tool
output against reference files.

The implementation is ANSI C style and uses standard C library calls plus
math functions (\texttt{log}, \texttt{log10}, \texttt{cos}). It is portable
across common Unix-like and Windows toolchains supported by STL.

\section{Example code}

The demonstration program is \texttt{gain\_chk.c}. It implements command
line parsing, frame-based spectral analysis, optional VAD-conditioned
statistics, report generation, and threshold-based return codes.
